\section{Model Design}

\subsection{Variational Autoencoder}
Variational autoencoders (VAEs)~\citep{Kingma2013AutoEncodingVB} are widely adopted in modern large-scale image and video generation models~\citep{rombach2022high} to reduce the computation of the subsequent diffusion model and facilitate efficient training and inference. Typically, a variational auto-encoder is usually composed of an encoder and a decoder; the encoder compresses the raw redundant pixel information into a compact latent representation, while the decoder reconstructs the original input from these latent features. The quality of VAE reconstruction directly establishes an upper bound for the realism and clarity achievable by the generative process, whereas the distribution of latent representations significantly impacts the convergence behavior of subsequent Diffusion Transformers (DiT). 

\textbf{Temporally-Causal Compression.} Following MAGVIT \cite{yu2023language}, we adopt a temporally causal convolutional architecture for both the encoder and decoder, allowing joint spatial-temporal compression of images and videos within latent space. To be more specific, the model transforms the input data from the RGB pixel space with shape $(T'+1, H', W', 3)$ into a continuous latent representation with shape $(T+1, H, W, C)$, where $(t, h, w, c)$ denotes time, height, width and channel dimensions with $r_t = \frac{T'}{T}$, $r_h = \frac{H'}{H}$, and $r_w = \frac{W'}{W}$ representing the downsample ratios along these three axes, respectively. Benefiting from the causal design, the VAE model can seamlessly process image input and output in the case of $T=T'=0$. The overall compression ratio is given by
\begin{equation}
    r = \frac{C \times T \times H \times W}{3 \times T' \times H' \times W'} = \frac{C}{3 \times r_t \times r_h \times r_w}.
    \label{equa:compression_ratio}
\end{equation}
In our practice, for the sake of training and inference efficiency and overall reconstruction and generation performance, we set $(r_t, r_h, r_w) = (4,16,16)$ and $C=48$. To accommodate the higher downsampling rate and pursue better generation performance, we remove the patchification operation on the DiT side, following the strategy adopted in DCAE \cite{chen2024deep}.

\textbf{VAE Training.} Our VAE is trained with $L1$ reconstruction loss, KL loss, LPIPS \cite{lpips_zhang2018unreasonable} perceptual loss and adversarial training loss.  Adversarial training has shown to be effective in improving the quality of VAE reconstruction by enforcing finer supervision on local textures and detailed structures. Taking into account appearance and motion modeling simultaneously, we apply a hybrid discriminator with an architecture similar to that used in PatchGAN \cite{isola2018imagetoimagetranslationconditionaladversarial}. 


\subsection{Diffusion Transformer}
With the visual tokens encoded by VAE and text tokens generated by a text encoder, we employ the transformer as our diffusion backbone~\cite{peebles2023scalable}, where a fine-tuned decoder-only LLM as the text encoder. The visual tokens are then concatenated with textual tokens and fed into the transformer blocks.


\textbf{Decoupled Spatial and Temporal Layers.} Considering both training and inference efficiency, we build the diffusion transformer with decoupled spatial and temporal layers, where the spatial layers perform attention aggregation within each frame, while the temporal layers focus attention computation across frames. We perform window partition within each frame in the temporal layers, allowing for a global receptive field across the temporal dimension. In addition, textual tokens only participate in cross-modality interaction in spatial layers. 


\begin{figure*}[t]
\centering
\includegraphics[width=\textwidth]{figures/framework.png}
\caption{Our diffusion transformer architecture.}
\end{figure*}

\textbf{MMDiT Architecture.} For the transformer blocks, we follow the MMDiT design in Stable Diffusion 3~\cite{esser2024scaling}, where a multi-modality self-attention layer is applied exclusively in spatial layers to integrate both the visual and textual tokens, whereas a self-attention layer only processes the visual tokens in temporal layers. Considering the semantic differences between visual and textual tokens, we use two separate sets of weights including adaptive layer norm, QKV projection, and MLP, for the two modalities in spatial layers. To prevent training instability, the Q and K embeddings are normalized prior to the attention matrix calculation.

\textbf{Multishot MM-RoPE.} In this paper, in addition to using 3D RoPE encoding for visual tokens, following Seaweed~\cite{seawead2025seaweed} and LCT~\cite{guo2025long}, we add 3D Multi-modal RoPE~(MM-RoPE) in the concatenated sequences by adding extra 1D positional encoding for textual tokens. The MM-RoPE also supports interleaved sequences of visual tokens and textual tokens, and can be extended to training video with multiple shots, where shots are organized in the temporal order of actions and each shot has its own detailed caption.

\textbf{Unified Task Formulation.} To enable conditional video generation, we concatenate the noisy inputs with cleaned or zero-padded frames along the channel dimension, and use binary masks to indicate which frames are instructions to follow~\cite{girdhar2024factorizing}. With this formulation, we can further unify different generation tasks such as text-to-image, text-to-video and image-to-video \cite{chen2023control}. During the training process, we mix these tasks and adjust the proportion by controlling the conditional inputs.


\subsection{Diffusion Refiner}
Take into account the training and inference efficiency, we employ a cascaded diffusion framework for high-resolution (HR) video generation. The base model generates 480p videos first,  which are then upscaled to 720p or 1080p high-resolution videos through a learned diffusion refiner model to enhance visual details and textures. 

\textbf{Refiner Model Training.} To facilitate training, the diffusion refiner model is initialized from the pre-trained base model. Different from the base model, the diffusion refiner model is trained with conditioning on the low-resolution 
(LR) videos. Specifically, the LR video is upsampled to a high resolution first, then concatenated with the diffusion noise along the channel dimension to form the input of the diffusion transformer. 


\subsection{Prompt Engineering (PE)}
\label{sec:pe}
As described in Sec~\ref{sec:caption}, texts used in DiT are form of dense video captions. Therefore, we need to employ a large language model to convert the user prompts into corresponding caption format. To achieve this, we initialize based on Qwen2.5-14B~\cite{yang2024qwen2_5} and employ two stages to implement high-quality Prompt Engineering (PE): Supervised Fine-Tuning (SFT) and Reinforcement Learning (RL).

\textbf{Supervised Fine-Tuning.} In the SFT stage, we synthesize large amount of user prompts and their dense caption expression by manual annotation. We specially devide the image-to-video (i2v) and text-to-video (t2v) tasks, as they are different in user prompt styles. We then adpot a fully fine-tuning strategy to train the model on the annotated data to aquire basic rephrasing abilitity.

\textbf{Reinforcement Learning.} However, due to the presence of model hallucinations, the results of the first SFT stage cannot guarantee that the semantics of the rewritten results fully meet the requirements of the user prompts. Therefore, we carefully collect a dataset of pairs with correct and incorrect rephrased results to perform the Direct Preference Optimization (DPO)~\cite{rafailov2023dpo,ji2024prompt} training. In this stage, we used the Low-Rank Adaptation (LoRA)~\cite{hu2022lora} fine-tuning strategy on the SFT model.

After the above stages, our prompt engineering model has strong ability to understand user prompts and gives precise and high-quality rephrased results in video caption format, consistent with DiT training.

